%% ---------------------------------------------------------------------
%% CHAPTER TEMPLATE  (copy to Overleaf/chapters/chNN-slug.tex and fill in)
%% Every chapter of the book follows this order.  Keep the section
%% headings; rename them to fit the topic, but do not drop a stage.
%% ---------------------------------------------------------------------
\chapter{Method Name}
\label{ch:chNN}
% Function names for pseudocode must be unique across the book: prefix them
% with a chapter-specific word.
\SetKwFunction{AlgMain}{MethodName}
\SetKwFunction{AlgHelper}{HelperName}

\begin{objectives}
\begin{itemize}
  \item Explain in one paragraph what the method computes and when it is the right tool.
  \item Trace the method by hand on a small instance.
  \item State the conditions under which it is correct / optimal, and why.
  \item Implement it and test it on a small example.
  \item Recognize the failure modes that matter in practice.
\end{itemize}
\end{objectives}

% ---------------------------------------------------------------------
\section{Why this matters}
\label{sec:chNN-motivation}
% 1/2 to 1 page.  A concrete scenario from the book's application in which this
% method is needed; where it sits in the system the book builds toward; a short
% roadmap of the chapter.

% ---------------------------------------------------------------------
\section{The idea in plain words}
\label{sec:chNN-intuition}
% Intuition before formalism.  A picture (\inputfigure{chNN/idea}) and
\begin{keyidea}
One or two sentences that a reader should remember for the rest of their life.
\end{keyidea}

% ---------------------------------------------------------------------
\section{Problem statement and notation}
\label{sec:chNN-problem}
% Formal definitions.  Use the macros of textbook.sty.
\begin{definition}[Name]
\label{def:chNN-name}
...
\end{definition}

% ---------------------------------------------------------------------
\section{The method}
\label{sec:chNN-algorithm}
% Pseudocode + line-by-line walkthrough (refer to lines with \ref{alg:...:line}).
\begin{algorithm}[htb]
\caption{Method name}
\label{alg:chNN-name}
\KwIn{...}
\KwOut{...}
\Fn{\AlgMain{$x$}}{
  initialize\label{alg:chNN-name:init}\;
  \While{not finished}{
    step\label{alg:chNN-name:step}\;
  }
  \KwRet result\;
}
\end{algorithm}

% ---------------------------------------------------------------------
\section{A worked example}
\label{sec:chNN-example}
% A small, fully worked instance with a figure and a trace table; every number
% comes from running code/chNN_slug.py.
\begin{example}[Descriptive title]
\label{ex:chNN-name}
...
\end{example}
\begin{figure}[tb]
  \centering
  \inputfigure{chNN/example}
  \caption{What the reader should notice in this figure.}
  \label{fig:chNN-example}
\end{figure}
\begin{table}[tb]
  \centering\small
  \caption{Trace of \cref{alg:chNN-name} on \cref{ex:chNN-name}.}
  \label{tab:chNN-trace}
  \begin{tabular}{@{}rlll@{}}
  \toprule
  Step & State & Value & Comment\\
  \midrule
  1 & ... & ... & ...\\
  \bottomrule
  \end{tabular}
\end{table}

% ---------------------------------------------------------------------
\section{Properties: correctness, guarantees, complexity}
\label{sec:chNN-properties}
\begin{theorem}[Name]
\label{thm:chNN-name}
...
\end{theorem}
\begin{proof}
... (a complete proof for the essential results; a proof sketch with a citation otherwise)
\end{proof}

% ---------------------------------------------------------------------
\section{Variants and extensions}
\label{sec:chNN-variants}

% ---------------------------------------------------------------------
\section{Implementation notes}
\label{sec:chNN-implementation}
% Data structures, numerical issues, then a listing excerpted verbatim from
% code/chNN_slug.py (ASCII only, at most 45 lines per listing).
\begin{lstlisting}[caption={Core of the implementation (from \code{code/chNN\_slug.py}).},label={lst:chNN-core}]
def method(x):
    ...
\end{lstlisting}

\begin{pitfall}
A common mistake, why it happens, and how to avoid it.
\end{pitfall}

% ---------------------------------------------------------------------
\section{Where this fits}
\label{sec:chNN-application}
\begin{applicationbox}
How the method is used in the system or practice the book builds toward, what it
receives from and delivers to the other parts, and where the course plan covers it.
\end{applicationbox}

% ---------------------------------------------------------------------
\section{Summary}
\begin{summary}
\begin{itemize}
  \item ...
\end{itemize}
\end{summary}

\paragraph{Further reading.}
Original source \cite{key}; textbook treatment \cite{key}.

% ---------------------------------------------------------------------
\section{Exercises}
\label{sec:chNN-exercises}
\begin{exercise}[\difficulty{1}]
\label{exr:chNN-slug}
...
\end{exercise}
\begin{exercise}[\difficulty{3} Coding]
\label{exr:chNN-coding}
The coding exercise the requirements assign to this topic.
\end{exercise}
